\chapter{Nacer humano}

\epigraph{Todos los mamíferos nacen. Pero solo el ser humano necesita que alguien le sostenga la cabeza.}{}

Hay algo que casi nadie se para a pensar, y que sin embargo dice mucho de lo que somos: el parto humano es, con diferencia, el más difícil del reino animal.

No hablo solo del dolor, aunque también. Hablo de algo más profundo. De un conflicto que lleva inscrito en nuestro cuerpo millones de años de historia. Un conflicto entre dos cosas que nos hacen humanos ---caminar erguidos y pensar--- y que, en el momento de nacer, chocan de frente. Para entender qué tiene esto que ver con el amor, con el matrimonio y con la familia, hay que empezar por el principio. Y el principio es la sabana africana, hace unos cuatro millones de años.

\section{Ponerse de pie}

Imaginemos a una hembra de \textit{Australopithecus afarensis} ---quizá deberíamos llamarla por su nombre más famoso: Lucy.\footnote{El esqueleto de Lucy fue descubierto en Etiopía en 1974 por Donald Johanson. Tiene unos 3,2 millones de años de antigüedad y fue clave para entender el origen del bipedalismo.} Lucy medía poco más de un metro, pesaba unos treinta kilos y hacía algo que ninguno de sus antepasados había hecho de forma habitual: caminaba sobre dos piernas.

No fue un capricho. El bipedalismo traía ventajas enormes. En una sabana cada vez más abierta, donde los bosques retrocedían, caminar erguido permitía ver por encima de la hierba alta ---detectar al depredador antes de que fuera demasiado tarde---. Liberaba las manos para cargar alimentos, para transportar herramientas rudimentarias, para sostener a una cría. Y era más eficiente: caminar sobre dos piernas gasta menos energía que hacerlo sobre cuatro.\footnote{Estudios comparativos muestran que el bipedalismo humano consume aproximadamente un 75\% menos de energía que la locomoción cuadrúpeda de un chimpancé de tamaño equivalente. Cf. Sockol, Raichlen y Pontzer, ``Chimpanzee locomotor energetics and the origin of human bipedalism'', \textit{PNAS}, 2007.}

Pero todo tiene un precio. Para caminar erguidos de manera eficiente, la pelvis tuvo que cambiar de forma. Se acortó, se ensanchó por arriba y se estrechó por abajo, formando una especie de cuenco compacto que sostenía las vísceras y transfería el peso a las piernas. La pelvis de Lucy ya no se parecía a la de un chimpancé. Se parecía cada vez más a la nuestra.

Y eso iba a ser un problema.

\section{El cerebro que no para de crecer}

Mientras la pelvis se estrechaba, algo más ocurría. Lentamente, generación tras generación, el cerebro de nuestros antepasados comenzó a crecer.

No fue solo un aumento de tamaño. Fue un cambio cualitativo. La corteza cerebral se plegaba en circunvoluciones cada vez más complejas. Un cambio de dieta tuvo mucho que ver: cuando nuestros ancestros empezaron a consumir más proteína animal ---primero carroña, luego caza--- y a cocinar los alimentos con fuego,\footnote{La ``hipótesis del cocinado'' de Richard Wrangham propone que la cocción de alimentos fue clave para el aumento del cerebro humano, al hacer los nutrientes más accesibles y reducir la energía necesaria para la digestión. Cf. Wrangham, \textit{Catching Fire: How Cooking Made Us Human}, Basic Books, 2009.} el intestino pudo acortarse y la energía sobrante se redirigió al órgano más caro del cuerpo. El cerebro humano supone apenas el 2\% del peso corporal, pero consume el 20\% de la energía.\footnote{En comparación, el cerebro de un primate típico consume alrededor del 8-10\% de la energía corporal.}

Más neuronas, más conexiones, más capacidad. Aquello nos permitió fabricar herramientas cada vez mejores, planificar, comunicarnos con un lenguaje articulado, cooperar en grupo. Nos estábamos convirtiendo en humanos.

Pero un cerebro más grande necesita una cabeza más grande para alojarlo. El cráneo de un \textit{Homo sapiens} recién nacido tiene un volumen medio de unos 350 centímetros cúbicos ---y seguirá creciendo después del parto hasta los 1.400 aproximados del adulto---. Compare eso con los 150-170 de un chimpancé neonato, que además nace con un grado de desarrollo mucho mayor.

Tenemos, entonces, dos fuerzas evolutivas que nos hacen humanos y que tiran en direcciones opuestas: una pelvis que se estrecha para caminar, y una cabeza que crece para pensar. El punto de encuentro de esas dos fuerzas es el canal del parto. Y ahí es donde estalla lo que los antropólogos llaman el \textbf{dilema obstétrico}.\footnote{El término ``dilema obstétrico'' fue acuñado por el antropólogo Sherwood Washburn en 1960. Aunque investigaciones recientes matizan algunos aspectos de la hipótesis original, el conflicto básico entre bipedalismo y encefalización sigue siendo ampliamente aceptado. Cf. Washburn, ``Tools and Human Evolution'', \textit{Scientific American}, 1960.}

\section{El dilema obstétrico}

El conflicto es real y medible. En la mayor parte de los primates, la cabeza del bebé pasa por el canal del parto con cierta holgura. En los humanos, no. La cabeza del bebé se ajusta al canal pélvico como una llave en su cerradura ---con apenas milímetros de margen---. Por eso el parto humano es más largo, más doloroso y más peligroso que el de cualquier otro primate. Por eso la cabeza del bebé tiene que rotar durante el descenso, ejecutando una especie de giro helicoidal para adaptarse a la forma cambiante del canal.\footnote{Este mecanismo de rotación es exclusivo del parto humano. En otros primates, el bebé desciende sin necesidad de rotar. Es también la razón por la que el parto humano se beneficia de la asistencia: la cría nace mirando hacia atrás, lo que dificulta que la propia madre la reciba. Somos una de las pocas especies en las que el parto es normalmente asistido.} Por eso los huesos del cráneo del recién nacido no están soldados: las fontanelas permiten que la cabeza se deforme ligeramente para pasar.

La evolución tenía que resolver este conflicto, y lo hizo con una solución ingeniosa pero de consecuencias enormes: adelantar el parto.

Conviene ser honesto: no todos los investigadores explican este adelanto de la misma manera. Hay otra hipótesis, propuesta por Holly Dunsworth y sus colegas, que apunta no tanto a la pelvis como al metabolismo.\footnote{Dunsworth, H. et al., ``Metabolic hypothesis for human altriciality'', \textit{PNAS}, 2012. Dunsworth argumenta que el parto se desencadena cuando la madre alcanza el límite de su capacidad metabólica ---unas 2 a 2,5 veces su tasa basal---, independientemente de las dimensiones pélvicas.} Según esta visión, lo que limita la duración del embarazo no es que el bebé no quepa, sino que la madre no puede sostener energéticamente un feto cuyo cerebro crece a una velocidad sin precedentes en el reino animal. El embarazo se interrumpe cuando el coste metabólico se vuelve insostenible.

¿Quién tiene razón? Probablemente ambas hipótesis capturan parte de la verdad. Pero lo que importa para nuestro argumento es lo que las dos comparten: en ambos casos, la causa última es el cerebro humano ---demasiado grande para la pelvis o demasiado costoso para el metabolismo---, y en ambos casos el resultado es el mismo: nacemos antes de estar listos.

\section{El feto externo}

Si un bebé humano naciera con el mismo grado de madurez neurológica que un chimpancé recién nacido, la gestación debería durar entre 18 y 21 meses.\footnote{Cf. Portmann, A., ``Die Ontogenese des Menschen als Problem der Evolutionsforschung'', \textit{Verhandlungen der Schweizerischen Naturforschenden Gesellschaft}, 1941. Portmann fue el primero en calcular esta gestación teórica y en proponer el concepto de altricialidad secundaria.} Pero con 18 meses de desarrollo ---ya sea por las dimensiones de la cabeza o por el coste energético---, nuestro cuerpo no puede más. La solución: nacer a los nueve meses, a medio camino.

El resultado es que los bebés humanos llegan al mundo radicalmente inmaduros. Los biólogos lo llaman \textit{altricialidad secundaria}: una especie que, por su linaje, debería producir crías relativamente maduras al nacer (como los demás primates), pero que en la práctica produce crías tan indefensas como las de un ratón o un gorrión.

Pensemos en lo que esto significa. Un potro se pone de pie y camina a las pocas horas de nacer. Un cervatillo puede seguir a su madre por el bosque en su primer día. Incluso un chimpancé recién nacido tiene la fuerza suficiente para agarrarse al pelaje de su madre.

Un bebé humano no puede hacer nada de eso. No puede sostener la cabeza. No puede agarrarse. No puede desplazarse. No regula bien su temperatura. Su visión es borrosa. Su sistema inmunitario es inmaduro. Durante meses, es absolutamente incapaz de sobrevivir sin que alguien le cuide de forma constante.

El antropólogo Ashley Montagu lo expresó con una imagen que se ha hecho célebre: el primer año de vida del ser humano es una \textit{exterogestación}, una gestación fuera del útero.\footnote{Montagu, A., \textit{Touching: The Human Significance of the Skin}, Harper \& Row, 1971.} El bebé ha nacido, sí, pero su desarrollo está lejos de completarse. Su cerebro triplicará su tamaño en los primeros dos años, formando las conexiones que lo harán pensar, hablar, sentir.

Nacemos, en cierto modo, a medio hacer. Y esa es la clave de todo lo que viene después.

\section{Una especie que necesita amor para sobrevivir}

Pensemos ahora en lo que significaba esto en la sabana africana hace doscientos mil años. Una madre acaba de dar a luz. Tiene en brazos a una criatura que depende absolutamente de ella. Que necesita ser alimentada cada pocas horas. Que no puede quedarse sola ni un momento. Que hay que cargar a todas partes, porque no va a caminar en un año.

Esa madre necesita comer, necesita beber, necesita dormir. Necesita estar alerta a los depredadores. ¿Cómo puede hacer todo eso con un bebé en brazos?

La respuesta es que no puede. No sola.

Y aquí es donde la historia da un giro que lo cambia todo. Porque la evolución no solo adelantó el parto. También preparó algo para lo que venía después: mecanismos biológicos que empujan a cuidar.

Durante el parto y la lactancia, el cuerpo de la madre se inunda de oxitocina.\footnote{La oxitocina es un neuropéptido producido en el hipotálamo. Además de su papel en las contracciones uterinas y la eyección de leche, actúa en el cerebro generando sensaciones de confianza, calma y vinculación. Se la ha llamado ``hormona del vínculo'', aunque su función es más compleja: en determinados contextos también puede intensificar la desconfianza hacia los extraños, reforzando la protección del grupo cercano.} No es una hormona cualquiera. Es la hormona que genera vínculo, que empuja a cuidar, que hace que una madre sienta que daría la vida por esa criatura que acaba de conocer. El contacto piel con piel dispara más oxitocina. La lactancia la mantiene alta. El llanto del bebé activa en el cerebro de la madre circuitos de alarma y respuesta que son anteriores a cualquier decisión consciente.\footnote{Estudios de neuroimagen muestran que el llanto de un bebé activa la amígdala y el córtex cingulado anterior de la madre en milisegundos, antes de que haya procesamiento consciente. Cf. Swain et al., ``Brain basis of early parent-infant interactions'', \textit{Journal of Child Psychology and Psychiatry}, 2007.}

No es que la madre ``decida'' amar a su hijo. Es que su cuerpo la empuja a ello con una fuerza que precede a la voluntad. La biología no sustituye al amor libre ---una madre siempre puede rechazar a su cría, y el amor verdadero será siempre un acto de libertad---, pero le pone el viento a favor. Le da un empujón poderoso en la dirección correcta.

Y no solo a la madre. También al padre.

\section{El padre que se queda}

En la mayoría de los mamíferos, el macho se desentiende de la cría tras la cópula. Cumple su función reproductiva y desaparece. Pero en los humanos, algo distinto ocurrió.

Los estudios muestran que cuando un hombre se convierte en padre, su biología cambia.\footnote{Cf. Gettler et al., ``Longitudinal evidence that fatherhood decreases testosterone in human males'', \textit{PNAS}, 2011. Este estudio, realizado con más de 600 hombres filipinos, mostró que los niveles de testosterona disminuían significativamente al convertirse en padres, y más aún en quienes se implicaban activamente en el cuidado.} Los niveles de testosterona bajan ---la hormona que impulsa la competición y la búsqueda de nuevas parejas--- y suben los de oxitocina y vasopresina ---hormonas vinculadas al cuidado y la protección---. El cerebro del padre también se reorganiza: se activan circuitos de empatía y respuesta al llanto del bebé que antes estaban más silenciosos.

¿Por qué la evolución favoreció esto? Por una razón implacable: en un entorno donde los bebés son tan vulnerables y la madre está tan exigida, aquellos grupos donde el padre permanecía, donde ayudaba a alimentar y proteger, tenían más probabilidades de que sus crías sobrevivieran. Y esas crías crecían, se reproducían, y transmitían esa misma inclinación. No la inclinación a ser ``buenos'', en abstracto, sino la inclinación biológica a vincularse, a quedarse, a cuidar.

El amor de pareja, el deseo de estar juntos, la satisfacción de construir algo común, tienen una raíz que va más allá del romanticismo: son mecanismos que aseguran que una cría extremadamente vulnerable llegue a la edad adulta. El amor no es un adorno de la vida. Es lo que hizo posible que la vida continuara.

\section{La vulnerabilidad que nos abre}

Hay algo asombroso en todo esto, si nos paramos a pensarlo.

Nuestra grandeza como especie ---ese cerebro extraordinario que nos permite componer sinfonías, curar enfermedades, preguntarnos por el sentido de la vida--- tiene un precio: nacemos incompletos. Nacemos necesitados. Nacemos absolutamente dependientes del amor de otros.

Y esa dependencia no es un fallo de diseño. Es el diseño.

Porque es precisamente nuestra vulnerabilidad la que nos abre al otro. Si fuéramos autosuficientes al nacer, como una tortuga que rompe el cascarón y se arrastra sola hacia el mar, no necesitaríamos vínculos. No necesitaríamos familias. No necesitaríamos el abrazo de una madre ni la presencia de un padre. No necesitaríamos amor.

Pero no somos tortugas. Somos la especie que nace con la cabeza sin cerrar, que llora para llamar a alguien, que necesita años de cuidado antes de poder valerse por sí misma. Somos la especie que, para sobrevivir, tuvo que aprender a amar.

Y quizá eso diga algo profundo sobre lo que somos. Quizá nuestro cuerpo, con su pelvis estrecha y su cerebro enorme, con sus hormonas que empujan al vínculo y su llanto que activa el cuidado, nos esté diciendo algo que intuimos pero que a veces olvidamos: que no estamos hechos para vivir solos. Que el amor no es un añadido opcional a la existencia humana, sino su condición de posibilidad.

Esa es la primera lección que nos da el cuerpo. Y es solo el comienzo.

\paracaminar

\begin{enumerate}
  \item ¿Alguna vez os habéis planteado que el cuerpo humano está ``diseñado'' para la relación? ¿Qué os sugiere esa idea?

  \item Si habéis tenido la experiencia de tener un bebé en brazos ---propio o ajeno---, ¿qué sentisteis? ¿Reconocéis esa ``fuerza'' biológica que empuja a cuidar?

  \item Hablad juntos: ¿qué significa para vosotros que el amor no sea solo un sentimiento, sino algo inscrito en el cuerpo?
\end{enumerate}

\vinetacapitulo{ch01}
